Le contour de chaque cercle est discrétisé par un polygone régulier inscrit. La longueur d'un côté de ce polygone est donnée par

\begin{equation}
l = 2R\sin\left(\frac{\pi}{N}\right),
\end{equation}

où $R$ désigne le rayon du cercle et $N$ le nombre de sommets. Le périmètre du polygone s'écrit alors

\begin{equation}
P_N
=
2NR\sin\left(\frac{\pi}{N}\right).
\label{eq:perimetre_poly}
\end{equation}

Considérons à présent deux cercles concentriques de rayons $R$ et $2R$, discrétisés respectivement par des polygones réguliers comportant $N_1$ et $N_2$ sommets. Le rapport des périmètres obtenus après discrétisation vaut

\begin{equation}
\frac{P_1^{(N)}}{P_2^{(N)}}
=
\frac{N_1\sin\left(\pi/N_1\right)}
     {2N_2\sin\left(\pi/N_2\right)}.
\label{eq:ratio_perimetres}
\end{equation}

Le rapport exact des périmètres des deux cercles est quant à lui égal à

\begin{equation}
\frac{P_1}{P_2}
=
\frac{2\pi R}{4\pi R}
=
\frac{1}{2}.
\end{equation}

L'erreur relative introduite par cette approximation géométrique est donc

\begin{equation}
\varepsilon_P
=
\frac{
\dfrac{P_1}{P_2}
-
\dfrac{P_1^{(N)}}{P_2^{(N)}}
}{
\dfrac{P_1}{P_2}
}
=
1-
\frac{N_1\sin\left(\pi/N_1\right)}
     {N_2\sin\left(\pi/N_2\right)}.
\label{eq:err_ratio_perimetre}
\end{equation}
